% Generated draft: requires mathematical and source review.
\documentclass[11pt]{article}
\usepackage[margin=25mm]{geometry}
\usepackage{fontspec}
\setmainfont{Latin Modern Roman}
\usepackage{amsmath,amssymb,mathtools,unicode-math}
\setmathfont{Latin Modern Math}
\usepackage{enumitem,needspace,xurl,hyperref}
\setlength{\parindent}{0pt}
\setlength{\parskip}{5pt}
\setlength{\emergencystretch}{3em}
\newcommand{\N}{\mathbb{N}}
\newcommand{\Z}{\mathbb{Z}}
\newcommand{\Q}{\mathbb{Q}}
\newcommand{\R}{\mathbb{R}}
\newcommand{\C}{\mathbb{C}}
\newcommand{\F}{\mathbb{F}}
\newcommand{\A}{\mathbb{A}}
\newcommand{\PP}{\mathbb{P}}
\newcommand{\E}{\mathbb{E}}
\newcommand{\eps}{\varepsilon}
\newcommand{\dd}{\,\mathrm{d}}
\begin{document}
\section*{UnsolvedMath ID: 30004106}
\textbf{Source problem code:} catalogue entry
\subsection{Definitions and mathematical statement}
Under the positive-characteristic interpretation of the supplied modular formulation, let \(p\) be a prime, let \(G\) be a finite group, and let \(k\) be a field of characteristic \(p\). No splitting-field or algebraic-closure hypothesis, and no assumption that \(p\mid |G|\), is supplied.

Let \(a(G)\) denote the Green ring of finite-dimensional left \(kG\)-modules: its additive relations are \([M\oplus N]=[M]+[N]\), and its multiplication is \([M][N]=[M\otimes_k N]\), with diagonal \(G\)-action on the tensor product. Its unit is the class of the trivial module \(k\). Choose representatives \((M_i)_{i\in I}\) of the isomorphism classes of nonzero indecomposable finite-dimensional \(kG\)-modules. By the direct-sum decomposition of finite-dimensional modules, their classes form an additive basis of \(a(G)\). Set
\[
 a_{\mathbb C}(G)=\mathbb C\otimes_{\mathbb Z}a(G),\qquad
 \left\|\sum_{i\in I}a_i[M_i]\right\|=\sum_{i\in I}|a_i|\dim_k M_i,
\]
where only finitely many coefficients are nonzero. The indecomposable-basis interpretation is necessary to make this displayed norm unambiguous.

For a module \(M\), let \(M^\vee=\operatorname{Hom}_k(M,k)\), with \((g\cdot f)(m)=f(g^{-1}m)\). Interpret the involution induced by duality as the conjugate-linear map
\[
 \left(\sum_i a_i[M_i]\right)^*=\sum_i\overline{a_i}[M_i^\vee].
\]
Let \(\widehat a(G)\) be the completion in the displayed norm, with the continuous extensions of multiplication and this involution. In particular, \(\widehat a(G)\) is a unital complex Banach \(*\)-algebra.

For a unital complex Banach algebra \(A\), write
\[
 \sigma_A(y)=\{\lambda\in\mathbb C:y-\lambda 1_A\text{ is not invertible in }A\}.
\]
A Banach \(*\)-algebra \(A\) is called symmetric if \(\sigma_A(x^*x)\subseteq[0,\infty)\) for every \(x\in A\).

The question is whether, for every \(G\) and \(k\) as above,
\[
 \forall x\in\widehat a(G),\qquad
 \sigma_{\widehat a(G)}(x^*x)\subseteq[0,\infty).
\]
The source writes only ``characteristic \(p\)''; the positive-characteristic scope used here follows its modular description. Whether characteristic zero was also intended remains unresolved. This is a statement-only draft requiring human mathematical review.
\subsection{Short English statement}
For every finite group and every field in the supplied modular setting, form the complexified Green ring, complete it using the dimension-weighted norm on indecomposable-module coefficients, and use conjugate-linear module duality as its involution. Must every element \(x^*x\) in this completion have spectrum contained in \([0,\infty)\)? The supplied record labels the problem open; that status has not been independently verified.
\subsection{Sources}
\begin{itemize}
\item Local UnsolvedMath registry: supplied complete record with permanent ID \(30004106\), title \emph{Symmetry of Completed Finite-Group Representation Rings}, source problem code \texttt{catalogue entry}, in the Oberwolfach Reports Open Problems collection. No individual proposer is identified in the supplied record.
\item Supplied report citation: \emph{Representations of Finite Groups (2020)}, Oberwolfach Reports, DOI \texttt{10.4171/owr/2019/14}, pp.\ 841--895. \url{https://doi.org/10.4171/owr/2019/14}. The record separately gives proposed year 2019 and an \texttt{owr-year:2019} tag; these are retained without reconciling them with the citation's 2020 designation.
\item Supplied literature source, labeled ``Benson on modular representation rings'': \url{https://homepages.abdn.ac.uk/d.j.benson/pages/papers/b/benson/banach-book.pdf}.
\item Supplied literature source, labeled ``Oberwolfach representation theory report'': \url{https://doi.org/10.4171/owr/2019/14}.
\item Supplied literature assessment dated exactly \texttt{2026-08-22}: the question is recorded as open, with the assessment that the primary reference explicitly calls it unknown and that no later source located in that curation establishes symmetry in all cases. This is supplied dated curation and literature triage, not an independently verified assessment; time-sensitive status and cited claims require independent checking.
\item The supplied statement-status label is \texttt{corrected\_verified}; its verification note reports comparison against the official Oberwolfach report and preservation of its hypotheses, notation, and target. These are source-record claims, not verification performed for this draft. No cited source was visited in preparing this statement.
\end{itemize}
\end{document}
